\section{Diseño del Subsistema de Visión Artificial}

El subsistema de visión artificial concentra la funcionalidad central del
proyecto: transformar el material visual capturado en los accesos en información
estructurada y verificable. A diferencia de un detector aislado, este subsistema
se organiza como un \textit{pipeline} en el que la salida de cada etapa alimenta
a la siguiente, por lo que sus interfaces de datos se especifican de forma
explícita para garantizar la trazabilidad y la verificabilidad de cada
componente.

\subsection{Arquitectura de los pipelines de visión}

El subsistema de visión comprende dos \textit{pipelines} con propósitos y
regímenes de ejecución distintos, cuya separación es deliberada.

El primero es el \textbf{pipeline de operación}, que se ejecuta en tiempo real
sobre el flujo de video de cada punto de acceso y conserva un estado por cámara
mientras un vehículo permanece dentro del campo de visión. Es el responsable de
producir los eventos de acceso, y es el que se describe en esta sección.

El segundo es el \textbf{pipeline de construcción del conjunto de datos}, que se
ejecuta fuera de línea sobre el material de video archivado. Su propósito no es
operar el parqueadero, sino producir las muestras etiquetadas con las que se
entrenan y evalúan los modelos; constituye, por tanto, un instrumento
metodológico, y se documenta junto con la construcción efectiva del conjunto de
datos en el Capítulo III.

Ambos comparten los módulos de detección y de reconocimiento, pero difieren en su
régimen de ejecución y en su producto: el primero emite eventos de acceso, el
segundo emite muestras anotadas. La Figura \ref{fig:pipeline_vision} presenta el
pipeline de operación, cuyas etapas y artefactos de datos se detallan en la Tabla
\ref{tab:modulos_vision}.

El pipeline de operación se organiza en seis etapas secuenciales. Dado que la
determinación de eventos depende de la trayectoria del vehículo, estas etapas se
ejecutan sobre una secuencia de fotogramas y no sobre capturas aisladas: el
estado de cada vehículo se conserva entre fotogramas dentro del alcance de su
cámara, lo que habilita el seguimiento y la interpretación de la trayectoria. El
servicio admite además un modo de captura puntual, en el que se recibe una única
imagen del vehículo; ese modo cubre las etapas de detección, reconocimiento de
placa y clasificación de atributos, pero no la determinación de eventos por
trayectoria, que requiere continuidad temporal.

\begin{figure}[htbp]
    \centering
    \begin{tikzpicture}
        % *** Entrada del pipeline ***
        \spData{entrada}{-2.6,1.15}{imagen del acceso}

        % *** Etapas del pipeline ***
        \spBlock{etapa1}{-2.6,0}{7.4cm}{matblue}
            {Captura y preprocesamiento}{RTSP, ROI y enmascarado}
        \spBlock{etapa2}{-2.6,-1.35}{7.4cm}{matblue}
            {Detección de vehículos}{YOLO26}
        \spBlock{etapa3}{-2.6,-2.70}{7.4cm}{matblue}
            {Seguimiento}{ByteTrack}
        \spBlock{etapa4}{-2.6,-4.05}{7.4cm}{matblue}
            {Localización y reconocimiento de placa}{YOLO26n + CRNN}
        \spBlock{etapa5}{-2.6,-5.40}{7.4cm}{matblue}
            {Clasificación de atributos}{Tipo y color}
        \spBlock{etapa6}{-2.6,-6.75}{7.4cm}{matblue}
            {Determinación de eventos}{Línea virtual}

        % *** Artefactos de datos ***
        \spData{d1}{4.5,0}{fotograma normalizado}
        \spData{d2}{4.5,-1.35}{cajas y clases}
        \spData{d3}{4.5,-2.70}{trayectoria con ID}
        \spData{d4}{4.5,-4.05}{texto de matrícula}
        \spData{d5}{4.5,-5.40}{atributos del vehículo}
        \spData{d6}{4.5,-6.75}{evento de acceso}

        % *** Flujo secuencial ***
        \draw[spflow] (entrada.south) -- (etapa1.north);
        \draw[spflow] (etapa1.south) -- (etapa2.north);
        \draw[spflow] (etapa2.south) -- (etapa3.north);
        \draw[spflow] (etapa3.south) -- (etapa4.north);
        \draw[spflow] (etapa4.south) -- (etapa5.north);
        \draw[spflow] (etapa5.south) -- (etapa6.north);

        % *** Artefactos emitidos por cada etapa ***
        \draw[spflow] (etapa1.east) -- (d1.west);
        \draw[spflow] (etapa2.east) -- (d2.west);
        \draw[spflow] (etapa3.east) -- (d3.west);
        \draw[spflow] (etapa4.east) -- (d4.west);
        \draw[spflow] (etapa5.east) -- (d5.west);
        \draw[spflow] (etapa6.east) -- (d6.west);

    \end{tikzpicture}
    \caption{Pipeline de operación del subsistema de visión artificial.}
    \label{fig:pipeline_vision}
\end{figure}

\begin{longtable}{
        P{0.18}
        P{0.20}
        P{0.20}
        P{0.32}
    }
    \caption{Módulos del pipeline de percepción}
    \label{tab:modulos_vision}\\

    \toprule
    \rowcolor{gray!20}
    \textbf{Módulo} &
    \textbf{Entrada} &
    \textbf{Salida} &
    \textbf{Restricción de diseño} \\
    \midrule
    \endfirsthead

    \multicolumn{4}{c}{\textit{Continuación de la Tabla
    \ref{tab:modulos_vision}}} \\
    \toprule
    \rowcolor{gray!20}
    \textbf{Módulo} &
    \textbf{Entrada} &
    \textbf{Salida} &
    \textbf{Restricción de diseño} \\
    \midrule
    \endhead

    Captura y preprocesamiento &
    Imagen del acceso &
    Fotograma normalizado &
    Debe recortar la región de interés y enmascarar la información incrustada
    por el grabador.
    \\

    \midrule
    Detección de vehículos &
    Fotograma normalizado &
    Cajas con clase y confianza &
    Debe operar sin post-procesamiento dependiente del número de objetos para
    mantener latencia constante.
    \\

    \midrule
    Seguimiento &
    Cajas detectadas &
    Trayectoria con identificador &
    Debe sobrevivir a oclusiones breves sin perder la identidad del vehículo.
    \\

    \midrule
    Localización y reconocimiento de placa &
    Recorte del vehículo &
    Cadena de matrícula &
    Debe localizar la placa dentro del vehículo y transcribirla sin requerir
    segmentación individual de caracteres ni depender de un umbral de
    binarización fijo.
    \\

    \midrule
    Clasificación de atributos &
    Recorte del vehículo &
    Tipo y color &
    Debe ser robusto frente a reflejos especulares y variaciones de iluminación.
    \\

    \midrule
    Determinación de eventos &
    Trayectoria del vehículo &
    Evento de ingreso o salida &
    Debe evitar disparos múltiples para un mismo vehículo.
    \\*
    \multicolumn{4}{r}{\footnotesize\textit{Fuente: Elaboración propia}} \\*
    \bottomrule
\end{longtable}

\subsection{Detección de vehículos}

Este módulo constituye la puerta de entrada analítica del sistema. Recibe el
fotograma normalizado y emite el conjunto de cajas delimitadoras con su clase y
su nivel de confianza, distinguiendo entre la silueta completa del vehículo y la
región de la placa de matrícula. La decisión de diseño central es el empleo de
un detector de una sola etapa con inferencia libre de supresión no máxima, lo
que garantiza que el tiempo de cómputo no dependa del número de objetos
presentes en la escena. Esta propiedad resulta indispensable en un acceso
vehicular, donde dos o tres vehículos pueden solaparse en el encuadre y la
supresión no máxima convencional podría descartar una detección válida o
introducir una latencia variable, tal como se fundamentó en el capítulo
anterior.

La salida de este módulo se organiza como una lista de detecciones, cada una
compuesta por la clase, las coordenadas de la caja, la confianza asociada y un
identificador de fotograma, lo que permite reconstruir posteriormente el origen
de cualquier decisión del sistema.

\subsection{Seguimiento vehicular}

El módulo de seguimiento asigna y mantiene un identificador único por vehículo a
lo largo del tiempo. Su función es impedir que un mismo automóvil sea tratado
como una sucesión de detecciones independientes, situación que generaría
registros duplicados de un único evento de acceso. Para ello, el módulo modela
el movimiento de cada caja mediante un estimador de estado que predice su
posición en el fotograma siguiente y la corrige con cada nueva detección,
combinando dicha predicción con la similitud geométrica entre las cajas.

La asociación entre predicciones y detecciones se resuelve priorizando las
coincidencias de alta confianza y recurriendo a las de confianza menor solo para
completar las trayectorias que permanecen sin asignar. Este criterio, propio del
enfoque de seguimiento en dos etapas, reduce las fragmentaciones de identidad
cuando un vehículo atraviesa oclusiones parciales provocadas por peatones, por
la propia barrera o por destellos de luz. Las trayectorias que dejan de recibir
detecciones se conservan durante un número limitado de fotogramas antes de
descartarse, y aquellas con un número insuficiente de observaciones se
consideran ruido y se eliminan. La salida del módulo es una trayectoria con
identificador, compuesta por la secuencia temporal de las posiciones del
vehículo y los atributos reconocidos.

\subsection{Reconocimiento de matrículas}

El módulo de reconocimiento recibe el recorte del vehículo y devuelve la cadena
alfanumérica correspondiente. Comprende dos etapas: la localización de la placa
dentro del recorte y la transcripción de la región localizada. Su diseño adopta
un enfoque libre de segmentación: en lugar de aislar y clasificar cada carácter,
un modelo de secuencia ---una red convolucional recurrente con decodificación
CTC--- procesa la imagen completa de la placa y emite la transcripción de forma
directa. Esta decisión responde a las condiciones reales del entorno, donde la
suciedad, los marcos metálicos o las sombras degradan la separación entre
caracteres y provocan el colapso de los enfoques clásicos por segmentación.

El módulo incorpora, además, un criterio de calidad previo al reconocimiento:
los recortes cuya resolución o nitidez resultan insuficientes, o cuya posición
en el encuadre no corresponde a la franja esperada de la placa, se descartan
antes de invocar al motor de reconocimiento. Se prevé también una etapa de
validación sintáctica opcional que contrasta el resultado con la gramática formal
de las matrículas del contexto de estudio. Su efecto no se presume: se mide sobre
el conjunto de prueba, porque una regla que reescribe posiciones puede convertir
una lectura correcta en incorrecta cuando el formato admite letras dentro del
bloque numérico. La implementación reporta ese resultado y la decisión adoptada
(Sección~\ref{sec:resultados_ocr}). La salida del módulo comprende la matrícula
reconocida, su cadena original y un indicador de confianza.

\subsection{Clasificación de atributos del vehículo}

Una vez localizado el vehículo, el sistema recorta su región y la somete a un
módulo de clasificación encargado de determinar el tipo de carrocería y el color
predominante.
Este módulo no interviene en la decisión de apertura, pero aporta una capa de
verificación cruzada: si la matrícula reconocida corresponde a un vehículo
registrado con un tipo o un color distintos a los observados, el sistema puede
señalar una posible suplantación de matrícula.

Para reducir la sensibilidad a los reflejos y a la iluminación variable, el
recorte se transforma a un espacio cromático que desacopla la información de
color de la intensidad luminosa antes de ingresar al clasificador. El resultado
se expresa como la categoría más probable junto con su confianza, y se conserva
únicamente si dicha confianza supera un umbral configurable; en caso contrario,
el atributo se registra como no determinado, evitando incorporar información
poco fiable al registro del evento.

\subsection{Determinación de eventos y contrato de datos}

La etapa final interpreta la trayectoria para decidir si el vehículo ingresó o
salió del recinto. Para ello se define una línea virtual sobre la calzada frente
a la barrera, complementada con dos regiones de control que delimitan el espacio
exterior y el interior del recinto. El evento se genera cuando el centro de la
caja del vehículo cruza dicha línea, y el sentido se determina a partir del
vector de desplazamiento respecto de la normal de la línea y de la región de
origen y destino de la trayectoria. Este criterio garantiza que cada vehículo
produzca exactamente un evento por cruce, evitando las falsas activaciones que
produciría un disparo basado únicamente en la presencia dentro de una región.

El resultado del \textit{pipeline} constituye la interfaz formal entre el
subsistema de visión y la capa de persistencia: un evento de acceso, no una
sesión. La Tabla \ref{tab:contrato_evento} define sus campos. La distinción
importa porque la sesión ---con su identificador, su estado y su duración--- no
viaja en el mensaje: se deriva al persistir el evento, que es lo único que la
visión puede afirmar con evidencia.

\begin{longtable}{
        Q{0.40}
        P{0.13}
        P{0.39}
    }
    \caption{Contrato de datos del evento de acceso}
    \label{tab:contrato_evento}\\

    \toprule
    \rowcolor{gray!20}
    \textbf{Campo} &
    \textbf{Tipo} &
    \textbf{Descripción} \\
    \midrule
    \endfirsthead

    \multicolumn{3}{c}{\textit{Continuación de la Tabla
    \ref{tab:contrato_evento}}} \\
    \toprule
    \rowcolor{gray!20}
    \textbf{Campo} &
    \textbf{Tipo} &
    \textbf{Descripción} \\
    \midrule
    \endhead

    \texttt{camera\_id} & texto &
    Identificador del punto de acceso que originó la observación.
    \\

    \midrule
    \texttt{track\_id} & entero &
    Identificador de la trayectoria asignada por el agregador; distingue dos
    vehículos de la misma cámara.
    \\

    \midrule
    \texttt{box} & entero[4] &
    Coordenadas en píxeles de la caja del vehículo, tomadas de la vista de mayor
    área de la trayectoria.
    \\

    \midrule
    \texttt{observed\_at} & marca temporal &
    Instante del fotograma en que se confirmó la trayectoria.
    \\

    \midrule
    \texttt{attributes}.\texttt{vehicle\_type} & texto &
    Categoría de carrocería identificada, o nulo si la confianza resulta
    insuficiente.
    \\

    \midrule
    \texttt{attributes}.\texttt{vehicle\_type\_confidence} & real &
    Confianza del clasificador de tipo; se conserva junto a la etiqueta para que
    las etapas posteriores puedan descartar lecturas débiles.
    \\

    \midrule
    \texttt{attributes}.\texttt{color} & texto &
    Color predominante identificado, o nulo si la confianza resulta insuficiente.
    \\

    \midrule
    \texttt{attributes}.\texttt{color\_confidence} & real &
    Confianza del clasificador de color.
    \\

    \midrule
    \texttt{detection\_confidence} & real &
    Confianza del detector de vehículos para la caja seleccionada.
    \\

    \midrule
    \texttt{plate}.\texttt{text} & texto &
    Cadena resultante del reconocimiento y de la validación sintáctica del
    formato de matrícula; ausente cuando no se leyó placa.
    \\

    \midrule
    \texttt{plate}.\texttt{confidence} & real &
    Confianza del reconocedor de caracteres, empleada por la política de
    emparejamiento.
    \\

    \midrule
    \texttt{gate} & enumeración &
    Punto de acceso de la observación: \texttt{entry} o \texttt{exit}.
    \\

    \midrule
    \texttt{image\_path} & texto &
    Ruta del recorte del vehículo conservado como evidencia; puede ser nulo.
    \\

    \midrule
    \texttt{model\_name}, \texttt{latency\_ms} & texto, real &
    Identificador del modelo empleado y tiempo de inferencia del fotograma.
    \\*
    \multicolumn{3}{r}{\footnotesize\textit{Fuente: Elaboración propia}} \\*
    \bottomrule
\end{longtable}
